\subsection{Estremi e separazione}
\label{sec:esercizi-fondamenti-estremi}

\begin{exercise}[Estremi di unione di intervalli]
\label{ex:fondamenti-estremi-unione-intervalli}
\difficoltaesercizio{0}
Per $A=[-1,2)\cup\{3\}$, determinare maggioranti e minoranti,
$\sup A$, $\inf A$, massimo e minimo, se esistono.
\rimandoteoria{def:maggiorante,def:minorante,def:massimo-minimo}
\end{exercise}
\begin{proof}[Soluzione dell'esercizio]
Poiché $-1\le x\le3$ per ogni $x\in A$ e $-1,3\in A$,
l'insieme dei maggioranti è $[3,+\infty)$ e quello dei
minoranti è $(-\infty,-1]$. Pertanto
\[
  \sup A=\max A=3,\qquad \inf A=\min A=-1.
\]
\end{proof}

\begin{exercise}[Estremi di un insieme indicizzato]
\label{ex:fondamenti-estremi-insieme-indicizzato}
\difficoltaesercizio{1}
Sia
\[
  A=\left\{1-\frac1n:n\in\N,\ n\ge1\right\}.
\]
Determinare $\sup A$, $\inf A$, massimo e minimo.
Giustificare l'assenza di un eventuale estremo mediante
la caratterizzazione con $\varepsilon$.
\rimandoteoria{thm:caratterizzazione-sup,def:massimo-minimo}
\end{exercise}
\begin{proof}[Soluzione dell'esercizio]
Per ogni $n\ge1$ si ha
\[
  0\le1-\frac1n<1,
\]
con uguaglianza a sinistra per $n=1$. Quindi
$\inf A=\min A=0$.

Per ogni $\varepsilon>0$, per la proprietà archimedea
esiste $n>1/\varepsilon$ e dunque
\[
  1-\varepsilon<1-\frac1n<1.
\]
Per la caratterizzazione dell'estremo superiore,
$\sup A=1$. Poiché $1\notin A$, il massimo non esiste.
\end{proof}

\begin{exercise}[Separazione di due insiemi]
\label{ex:fondamenti-separazione-due-insiemi}
\difficoltaesercizio{1}
Siano $A=(-2,1)$ e $B=[1,5]$. Verificare che sono separati,
determinare tutti i loro elementi separatori e calcolare
$\sup A$ e $\inf B$.
\rimandoteoria{sec:lez06-insiemi-separati}
\end{exercise}
\begin{proof}[Soluzione dell'esercizio]
Per ogni $a\in A$ e $b\in B$ si ha $a<1\le b$,
quindi i due insiemi sono separati. Inoltre,
\[
  \sup A=1,\qquad \inf B=1.
\]
Un separatore $\lambda\in\R$ deve soddisfare
\[
  \sup A\le\lambda\le\inf B,
  \quad\text{ossia}\quad 1\le\lambda\le1.
\]
L'unico elemento separatore è pertanto $\lambda=1$.
\end{proof}

\begin{exercise}[Estremi di un'immagine razionale]
\label{ex:fondamenti-estremi-immagine-razionale}
\difficoltaesercizio{2}
Sia
\[
  A=\left\{\frac{x^2}{1+x^2}:x\in\R\right\}.
\]
Determinare l'immagine $A$, i suoi estremi superiore e inferiore
ed eventualmente massimo e minimo. Non utilizzare derivate.
\rimandoteoria{def:lez06-immagine,def:massimo-minimo}
\end{exercise}
\begin{proof}[Soluzione dell'esercizio]
Per ogni $x\in\R$, poiché $1+x^2>0$, si ha
\[
  0\le\frac{x^2}{1+x^2}
    =1-\frac1{1+x^2}<1.
\]
Viceversa, per ogni $y\in[0,1)$, scegliendo
\[
  x=\sqrt{\frac{y}{1-y}}\in\R
\]
si ottiene $x^2/(1+x^2)=y$. Pertanto
\[
  A=[0,1),\qquad \inf A=\min A=0,\qquad \sup A=1.
\]
Il massimo non esiste, poiché $1\notin A$.
\end{proof}

\begin{exercise}[Estremi di una famiglia alterna]
\label{ex:fondamenti-estremi-famiglia-alterna}
\difficoltaesercizio{2}
Si consideri
\[
  A=\left\{(-1)^n\left(1-\frac1n\right):
    n\in\N,\ n\ge1\right\}.
\]
Determinare $\sup A$, $\inf A$, massimo e minimo,
distinguendo gli indici pari da quelli dispari.
Dimostrare rigorosamente quali valori siano raggiunti.
\rimandoteoria{thm:caratterizzazione-sup,thm:caratterizzazione-inf}
\end{exercise}
\begin{proof}[Soluzione dell'esercizio]
Per $n$ pari, con $n\ge2$, si ha
\[
  \frac12\le1-\frac1n<1.
\]
Per $n$ dispari, con $n\ge1$, si ha
\[
  -1<-1+\frac1n\le0.
\]
Pertanto $-1<a<1$ per ogni $a\in A$.

Fissato $\varepsilon>0$, per la proprietà archimedea
si possono scegliere un indice pari $n>1/\varepsilon$
e un indice dispari $m>1/\varepsilon$. Allora
\[
  1-\varepsilon<1-\frac1n<1,
  \qquad
  -1<-1+\frac1m<-1+\varepsilon.
\]
Per la caratterizzazione degli estremi si conclude
\[
  \sup A=1,\qquad \inf A=-1.
\]
Poiché $-1,1\notin A$, non esistono né massimo né minimo.
\end{proof}

\begin{exercise}[Estremo superiore di un'unione]
\label{ex:fondamenti-sup-unione}
\difficoltaesercizio{2}
Siano $A,B\subseteq\R$ non vuoti e superiormente limitati.
Dimostrare che
\[
  \sup(A\cup B)=\max\{\sup A,\sup B\}.
\]
\rimandoteoria{def:estremo-superiore}
\end{exercise}
\begin{proof}[Soluzione dell'esercizio]
Si pone $M:=\max\{\sup A,\sup B\}$. Poiché
\[
  a\le\sup A\le M\quad\forall a\in A,
  \qquad
  b\le\sup B\le M\quad\forall b\in B,
\]
$M$ è un maggiorante di $A\cup B$.

Se $K$ è un maggiorante di $A\cup B$, allora è un
maggiorante sia di $A$ sia di $B$; pertanto
\[
  \sup A\le K,\qquad \sup B\le K,
  \qquad M\le K.
\]
Quindi $M$ è il minimo dei maggioranti di $A\cup B$,
ossia $\sup(A\cup B)=M$.
\end{proof}

\begin{exercise}[Estremo superiore di un'intersezione]
\label{ex:fondamenti-sup-intersezione}
\difficoltaesercizio{2}
Siano $A,B\subseteq\R$ non vuoti e superiormente limitati
e si supponga $A\cap B\ne\varnothing$. Dimostrare che
\[
  \sup(A\cap B)\le\min\{\sup A,\sup B\}.
\]
Costruire un esempio in cui la disuguaglianza è stretta.
\rimandoteoria{def:estremo-superiore}
\end{exercise}
\begin{proof}[Soluzione dell'esercizio]
Poiché $A\cap B\subseteq A$ e $A\cap B\subseteq B$,
per ogni $x\in A\cap B$ si ha
\[
  x\le\sup A,\qquad x\le\sup B,
\]
e quindi $x\le\min\{\sup A,\sup B\}$. Pertanto
\[
  \sup(A\cap B)\le\min\{\sup A,\sup B\}.
\]

Per esempio, scegliendo
\[
  A=[0,2]\cup\{5\},\qquad B=[0,2]\cup\{4\},
\]
si ha $A\cap B=[0,2]$ e dunque
\[
  \sup(A\cap B)=2<4=\min\{\sup A,\sup B\}.
\]
\end{proof}
